%% kapa: exam=1 round=20 year=2009 period=2 form=A subject=accounting subject_ko=회계학 seq=39 number=39
%% source: Q-Net 감정평가사 기출문제 (cst00309) artlSeq=5024178 file=extracted/1008/5024178/2009년 제20회 감정평가사 1차시험 기출문제/2교시.hwp
%% conversion: kordoc 4.24.51 -> tools/gen1.py (auto); review=proofread-sonnet; flags=none
\begin{question}{39}
다음은 (주)한국의 제품별 예산자료의 일부이다.

{\small\centering
\begin{tabular}{@{}l r r@{}}
\hline
 & \multicolumn{1}{c}{제품 A} & \multicolumn{1}{c}{제품 B} \\
\cline{2-2}\cline{3-3}
단위당 판매가격 & ￦400 & ￦500 \\
단위당 변동비 & ￦150 & ￦300 \\
단위당 기계시간 & 4시간 & 2시간 \\
최대 수요량(연간) & 100단위 & 200단위 \\
\hline
\end{tabular}\par}
사용가능한 총 기계시간이 연간 500시간일 때, 이익을 극대화하기 위해서는 두 제품을 각각 몇 단위씩 생산․판매하여야 하는가? \par

{\small\centering
\begin{tabular}{@{}ccc@{}}
 & 제품 A & 제품 B \\
\cline{2-2}\cline{3-3}
① & 25단위 & 150단위 \\
② & 25단위 & 200단위 \\
③ & 50단위 & 150단위 \\
④ & 50단위 & 200단위 \\
⑤ & 100단위 & 50단위 \\
\end{tabular}\par}
\end{question}
\answer{2}
